%======================================================================
\section{\texttt{exit}}
%======================================================================

\begin{mathpar}
  \inferrule* [lab=exit]
    {|S| \ge 1 \\ H(F.\rid) = R \\ R.\status = \Open}
    {\langle S :: F, H \rangle \xrightarrow{\texttt{exit}} \langle S,\ H[F.\rid \mapsto R'] \rangle
     \quad\text{where } R' = R \text{ with } \status := \Closed}
\end{mathpar}

\texttt{exit} pops the top frame $F$ and closes its region. Throughout this section fix
$\sigma = \langle S :: F, H\rangle$ with $\sigma \vdash \inv{ValidConfig}$, and
$\sigma' = \langle S, H'\rangle$ with $H' = H[F.\rid \mapsto R']$ as above. Let $n = |S|$, the
index of $F$ in $\sigma$. This is exactly what \texttt{exit\_cases} extracts from
$\texttt{exit}\ \sigma = \mathsf{some}\ \sigma'$. Lemma~\ref{lem:status-update} applies with
$rid_0 = F.\rid$ and $s = \Closed$.

The proof has three parts: one lemma that describes every location of $\sigma'$ in terms of $\sigma$
(Lemma~\ref{lem:exit-loc}), the observation that nothing new is stored (Lemma~\ref{lem:exit-refs}),
and the ten invariants.

\subsection{Supporting lemmas}

\begin{lemma}[Locations after \texttt{exit}]
\label{lem:exit-loc}
\lean{exit\_loc}
For every reference $r$, either $\loc{\sigma'}{r} = \loc{\sigma}{r}$, or $\loc{\sigma'}{r} = \noloc$
and $\loc{\sigma}{r} = \Stk(i)$ for some $i$.
\end{lemma}
\begin{proof}
For $r = \RId(rid)$: the location depends only on $\dom$, which is unchanged by
Lemma~\ref{lem:status-update}(a), so $\loc{\sigma'}{r} = \loc{\sigma}{r}$.

For $r = \OId(oid)$: by Lemma~\ref{lem:status-update}(d), the heap search gives the same answer in
$H'$ as in $H$, so only the stack search can differ. The frames of $\sigma'$ are those of $\sigma$
minus the top one, with the same indices. Searching $\sigma$'s stack from the top therefore first
looks at $(F, n)$, and if $oid \notin \dom(F.\omap)$ then continues exactly as the search of
$\sigma'$'s stack.
\begin{itemize}[nosep]
  \item If $oid \notin \dom(F.\omap)$, both searches return the same result, so the two locations
    agree.
  \item If $oid \in \dom(F.\omap)$, the stack search in $\sigma$ returns $(F, n)$. By \inv{L1},
    $oid$ is in no region of $H$, so the heap search fails and $\loc{\sigma}{r} = \Stk(n)$. Again by
    \inv{L1}, $oid$ is in no frame of $S$, since $\oids(S :: F) = \oids(S) \cat \dom(F.\omap)$ is
    duplicate-free. So in $\sigma'$ both searches fail, and $\loc{\sigma'}{r} = \noloc$. \qedhere
\end{itemize}
\end{proof}

\begin{corollary}
\label{cor:exit-loc}
For every reference $r$:
\begin{enumerate}[nosep,label=(\alph*)]
  \item $\loc{\sigma}{r} = \Rgn(rid)$ iff $\loc{\sigma'}{r} = \Rgn(rid)$;
    \leanpart{exit\_corollary\_1}
  \item if $\loc{\sigma'}{r} = \Stk(i)$ then $\loc{\sigma}{r} = \Stk(i)$.
    \leanpart{exit\_corollary\_3}
\end{enumerate}
\end{corollary}
\begin{proof}
In the second case of Lemma~\ref{lem:exit-loc}, $\loc{\sigma'}{r}$ is undefined and
$\loc{\sigma}{r}$ is a stack location, so neither side of (a) holds and the premise of (b) is false.
In the first case both statements are trivial.
\end{proof}

\begin{lemma}[No new references]
\label{lem:exit-refs}
\lean{exit\_refs\_mem}
Every element of $\refs(\sigma')$ is an element of $\refs(\sigma)$. More precisely,
$\refs(S) \sublist \refs(S :: F)$ and $\refs(H') \perm \refs(H)$.
\end{lemma}
\begin{proof}
$\refs(S :: F) = \refs(S) \cat \refs(F)$, and $\refs(H') \perm \refs(H)$ is
Lemma~\ref{lem:status-update}(b).
\end{proof}

\subsection{The invariants}

\begin{lemma}[\inv{L1}] \lean{exit\_L1}
$\sigma' \vdash \inv{L1}$.
\end{lemma}
\begin{proof}
$\oids(\sigma') = \oids(S) \cat \oids(H') \perm \oids(S) \cat \oids(H)$ by
Lemma~\ref{lem:status-update}(b), and this is a sublist of
$\oids(S) \cat \dom(F.\omap) \cat \oids(H) = \oids(\sigma)$, which is duplicate-free by \inv{L1}.
\end{proof}

\begin{lemma}[\inv{L2}] \lean{exit\_L2}
$\sigma' \vdash \inv{L2}$.
\end{lemma}
\begin{proof}
Let $G \in S$. By \inv{S1} for $\sigma$, the region ids of $S :: F$ are pairwise distinct, so
$G.\rid \ne F.\rid$. Hence $H'(G.\rid) = H(G.\rid)$, which by \inv{L2} for $\sigma$ (as
$G \in S :: F$) is an $\Open$ region.
\end{proof}

\begin{lemma}[\inv{H1}] \lean{exit\_H1}
$\sigma' \vdash \inv{H1}$.
\end{lemma}
\begin{proof}
A region of $H'$ is either $R'$ or a region of $H$ stored at some key other than $F.\rid$. In the
second case \inv{H1} for $\sigma$ applies directly. In the first, $R'$ has the same $\bridge$ and
$\omap$ as $R$, and $R$ satisfies \inv{H1} since it is a region of $H$.
\end{proof}

\begin{lemma}[\inv{H2}] \lean{exit\_H2}
$\sigma' \vdash \inv{H2}$.
\end{lemma}
\begin{proof}
Fix $rid$ and write $\#\ell$ for the number of occurrences of $\RId(rid)$ in $\ell$. By
Lemma~\ref{lem:exit-refs}, $\#\refs(S) \le \#\refs(S :: F)$ and $\#\refs(H') = \#\refs(H)$, so
$\#\refs(S) + \#\refs(H') \le \#\refs(S :: F) + \#\refs(H) \le 1$ by \inv{H2} for $\sigma$.
\end{proof}

\begin{lemma}[\inv{H3}] \lean{exit\_H3}
$\sigma' \vdash \inv{H3}$.
\end{lemma}
\begin{proof}
Let $H'(rid) = Q'$ and $\OId(oid) \in \refs(Q')$. By Lemma~\ref{lem:status-update}(c),
$H(rid) = Q$ for some $Q$ with the same $\omap$ as $Q'$, and $\refs$ of a region depends only on its
$\omap$, so $\OId(oid) \in \refs(Q)$. By \inv{H3} for $\sigma$,
$\loc{\sigma}{\OId(oid)} = \Rgn(rid)$, and by Corollary~\ref{cor:exit-loc}(a),
$\loc{\sigma'}{\OId(oid)} = \Rgn(rid)$.
\end{proof}

\begin{lemma}[\inv{S1}] \lean{exit\_S1}
$\sigma' \vdash \inv{S1}$.
\end{lemma}
\begin{proof}
$[G.\rid \mid G \in S]$ is a sublist of $[G.\rid \mid G \in S :: F]$, which is duplicate-free by
\inv{S1} for $\sigma$.
\end{proof}

\begin{lemma}[\inv{S2}] \lean{exit\_S2}
$\sigma' \vdash \inv{S2}$.
\end{lemma}
\begin{proof}
Let $(G, i) \in S$, $\OId(oid) \in \refs(G)$ and $\loc{\sigma'}{\OId(oid)} = \Stk(j)$. The pair
$(G,i)$ is also a frame of $\sigma$, at the same index. By Corollary~\ref{cor:exit-loc}(b),
$\loc{\sigma}{\OId(oid)} = \Stk(j)$, so \inv{S2} for $\sigma$ gives $j \le i$.
\end{proof}

\begin{lemma}[\inv{S3}] \lean{exit\_S3}
$\sigma' \vdash \inv{S3}$.
\end{lemma}
\begin{proof}
Let $(G, i) \in S$, $\OId(oid) \in \refs(G)$ and $\loc{\sigma'}{\OId(oid)} = \Rgn(rid')$. As in the
proof of \inv{S2}, $(G, i)$ is a frame of $\sigma$, and by Corollary~\ref{cor:exit-loc}(a)
$\loc{\sigma}{\OId(oid)} = \Rgn(rid')$. \inv{S3} for $\sigma$ gives a frame $(G'', k)$ of $\sigma$
with $G''.\rid = rid'$ and $k \le i$. Since $i < n$ we have $k < n$, so $(G'', k)$ is not the popped
frame: $(G'', k) \in S$, which is the required witness in $\sigma'$.
\end{proof}

\begin{lemma}[\inv{HS1}] \lean{exit\_HS1}
$\sigma' \vdash \inv{HS1}$.
\end{lemma}
\begin{proof}
Let $\OId(oid) \in \refs(\sigma')$. By Lemma~\ref{lem:exit-refs} it is in $\refs(\sigma)$, so by
\inv{HS1} for $\sigma$,
\[ oid \in \oids(\sigma) = \oids(S) \cat \dom(F.\omap) \cat \oids(H). \]
If $oid \in \oids(S)$ it is in $\oids(\sigma')$ directly; if $oid \in \oids(H)$ it is in
$\oids(H') \perm \oids(H)$ (Lemma~\ref{lem:status-update}(b)).

It remains to rule out $oid \in \dom(F.\omap)$: the popped frame's objects must not still be
referenced. Suppose $oid \in \dom(F.\omap)$. By Lemma~\ref{lem:loc-l1}(b),
$\loc{\sigma}{\OId(oid)} = \Stk(n)$. Now $\OId(oid)$ is in $\refs(\sigma') = \refs(S) \cat \refs(H')$:
\begin{itemize}[nosep]
  \item If $\OId(oid) \in \refs(G)$ for some $(G, m) \in S$, then $(G, m)$ is a frame of $\sigma$ and
    \inv{S2} for $\sigma$ gives $n \le m$, contradicting $m < n$.
  \item If $\OId(oid) \in \refs(H')$, then by Lemma~\ref{lem:status-update}(b) it is in $\refs(Q)$
    for some $H(rid) = Q$. \inv{H3} for $\sigma$ gives $\loc{\sigma}{\OId(oid)} = \Rgn(rid)$,
    contradicting $\loc{\sigma}{\OId(oid)} = \Stk(n)$. \qedhere
\end{itemize}
\end{proof}

\begin{lemma}[\inv{HS2}] \lean{exit\_HS2}
$\sigma' \vdash \inv{HS2}$.
\end{lemma}
\begin{proof}
If $\RId(rid) \in \refs(\sigma')$ then $\RId(rid) \in \refs(\sigma)$ by Lemma~\ref{lem:exit-refs}, so
$rid \in \dom(H) = \dom(H')$ by \inv{HS2} for $\sigma$ and Lemma~\ref{lem:status-update}(a).
\end{proof}

\begin{theorem} \lean{exit\_valid}
If $\sigma \vdash \inv{ValidConfig}$ and $\sigma \xrightarrow{\texttt{exit}} \sigma'$, then
$\sigma' \vdash \inv{ValidConfig}$.
\end{theorem}
\begin{proof}
The ten lemmas above.
\end{proof}

